\documentclass[14pt,aspectratio=1610]{beamer}
% \documentclass[12pt,aspectratio=1610,handout]{beamer}
\usepackage{csu-theme}
\usepackage{arrays}
\usepackage{linked-lists}
\usepackage{booktabs}
\usetikzlibrary{backgrounds}

\title{\Huge{} The Standard Template Library (STL)}
\subtitle{CSCI 315: Data Structure Analysis\\Chapter 21}
%\author{Sean T. Hayes, Ph.D.}
\institute{Charleston Southern University}

%% Comment out date to use default (today)
% \date{March 20, 2025}
\subject{Software Programming}
%\keywords{}
\graphicspath{{../imgs/stl}}

\setlength{\parskip}{0.5em}

\begin{document}

\begin{frame}
  \titlepage{}
\end{frame}

\section{Introduction}

\begin{frame}{Introduction to the STL}

\begin{itemize}
\item
  The standard C++ is equipped with a \emph{Standard Template Library\\(STL)}.
\item
  The STL includes class templates to process lists, stacks, and queues.
\item
  In this lecture, you will:
  \begin{itemize}
  \item
    Be introduced the Standard Template Library (STL)
  \item
    Become familiar with the three basic components of the STL:\@ containers, iterators, and algorithms.
  \item
    Become familiar with basic operations on vector objects
  \end{itemize}
\end{itemize}
\end{frame}

\begin{frame}{Components of the STL}
The STL consists of three categories of components:
\begin{enumerate}
\item
  \emph{Containers} are class templates to manage objects of a given type.
\item
  \emph{Iterators} are used to step through the elements of a container.
\item
  \emph{Algorithms} are used to manipulate data.
\end{enumerate}
\end{frame}

\begin{frame}{Components of the STL}
The STL consists of three categories of components:
\centerline{%
  \includegraphics[width=0.9\textwidth]{stl-componenets}%
}
\end{frame}

\begin{frame}{Available Containers Types}
\begin{itemize}[<+->]
\item
  \emph{Sequence Containers} implement data structures \\that can be accessed sequentially.
\item
  \emph{Ordered Associative Containers} implement sorted data structures that can be quickly searched (\(O(\log n)\) complexity).
\item
  \emph{Unordered Associative Containers} implement unsorted (hashed) data structures that can be quickly searched (\(O(1)\) on average, \(O(n)\) worst-case complexity).
\item
  \emph{Adapters} provide a different interface for sequential containers.
\item
  \emph{Views} provide flexible lightweight non-owning \\access to contiguous contains.
\end{itemize}
\end{frame}

\begin{frame}{Sequence Containers Types in C++20}
\renewcommand{\arraystretch}{1.1}
\begin{tabular}{l l l}
\toprule
  \textbf{Sequence Containers} & \textbf{Header File} & \textbf{Iterator Support}\\
  \midrule
  \href{https://en.cppreference.com/w/cpp/container/array}{array}
    & \lstinline{<array>} & Random access\\
  \href{https://en.cppreference.com/w/cpp/container/vector}{vector}
    & \lstinline{<vector>} & Random access\\
  \href{https://en.cppreference.com/w/cpp/container/deque}{deque}
    & \lstinline{<deque>} & Random access\\
    \href{https://en.cppreference.com/w/cpp/container/list}{list}
    & \lstinline{<list>} & Bidirectional\\
  \href{https://en.cppreference.com/w/cpp/container/forward_list}{forward\_list}
    & \lstinline{<forward_list>} & Forward\\
\bottomrule
\end{tabular}
\end{frame}


\begin{frame}{Associative Containers Types in C++20}
  \renewcommand{\arraystretch}{1.1}
  \begin{tabular}{l l l}
  \toprule
  \textbf{Associative Containers} & \textbf{Header File} & \textbf{Iterator Support}\\
  \midrule

    \href{https://en.cppreference.com/w/cpp/container/map}{map}
      & \lstinline{<map>} & Bidirectional\\
    \href{https://en.cppreference.com/w/cpp/container/multimap}{multimap}
      & \lstinline{<map>} & Bidirectional\\
    \href{https://en.cppreference.com/w/cpp/container/set}{set}
      & \lstinline{<set>} & Bidirectional\\
    \href{https://en.cppreference.com/w/cpp/container/multiset}{multiset}
      & \lstinline{<set>} & Bidirectional\\
  \midrule
    \href{https://en.cppreference.com/w/cpp/container/unordered_map}{unordered\_map}
      & \lstinline{<unordered\_map>} & Bidirectional\\
    \href{https://en.cppreference.com/w/cpp/container/unordered_multimap}{unordered\_multimap}
      & \lstinline{<unordered\_map>} & Bidirectional\\
    \href{https://en.cppreference.com/w/cpp/container/unordered_set}{unordered\_set}
      & \lstinline{<unordered\_set>} & Bidirectional\\
    \href{https://en.cppreference.com/w/cpp/container/unordered_multiset}{unordered\_multiset}
      & \lstinline{<unordered\_set>} & Bidirectional\\
  \bottomrule
\end{tabular}

\end{frame}


\begin{frame}{Adpaters and Views Types in C++20}
  \renewcommand{\arraystretch}{1.1}
  \begin{tabular}{l l l}
    \toprule
      \textbf{Adapters} & \textbf{Header File} & \textbf{Iterator Support}\\
    \midrule
      \href{https://en.cppreference.com/w/cpp/container/stack}{stack}
        & \lstinline{<stack>} & No iterator support\\
      \href{https://en.cppreference.com/w/cpp/container/queue}{queue}
        & \lstinline{<queue>} & No iterator support\\
      \href{https://en.cppreference.com/w/cpp/container/priority_queue}{priority queue}
        & \lstinline{<queue>} & No iterator support\\
    \midrule
      \textbf{Views} & \textbf{Header File} & \textbf{Iterator Support}\\
        \href{https://en.cppreference.com/w/cpp/container/span}{span}
        & \lstinline{<span>} & Bidirectional\\
  \bottomrule
  \end{tabular}
  
  \end{frame}

\begin{frame}{Sequence Containers in the STL}
Sequence container: every object in the container has a specific position.

\begin{enumerate}
\item
  \lstinline{array} -- encapsulates a fixed size array
\item
  \lstinline{vector} -- a list with elements stored contiguously
\item
  \lstinline{deque} -- a double-ended queue
\item
  \lstinline{list} --  usually a doubly-linked list
  \item
  \lstinline{forward\_list} --  usually a singly-linked list
\end{enumerate}
\end{frame}


\section{\texttt{\bfseries vector}}

\begin{frame}[fragile]{Declaring a \texttt{\bfseries vector}}
\begin{itemize}
\item
  Stores and manages its objects in a dynamic array
\item
  Must use: \lstinline{#include <vector>}
\item
  To define an object of type vector, specify the type of the object
\item
  Examples:
  \begin{lstlisting}
std::vector<int> intList;
std::vector<std::string> stringList;
\end{lstlisting}
\item
  \lstinline{vector} contains several constructors.
\end{itemize}
\end{frame}

\begin{frame}[fragile]{Constructors for a \texttt{\bfseries vector}}
  \small
There are 10 different constructors for \lstinline{vector}; Here are some highlights.
\begin{itemize}
\item
  \lstinline{vector();}\\
  Default constructor. Constructs an empty container.
\item
  \lstinline{vector(size_type count, const T& value = T());}\\
  Constructs the container with count copies of elements with value value.
\item
  \lstinline{template<typename InputIt>}\\
  \lstinline{vector(InputIt first, InputIt last);}\\
  Constructs the container with the contents of the range [first, last).
\item
  \lstinline{vector(const vector& other);}\\
  Copy constructor.
\end{itemize}
\end{frame}

\begin{frame}[fragile]{Operations on a \texttt{\bfseries vector}}
  \small
\begin{itemize}
\item
  Basic vector operations
  \begin{itemize}
  \item
    Item insertion\\
    \lstinline{insert(iterator pos, T& value);} \lstinline{push_back(const T& value);} etc.
  \item
    Item deletion\\
    \lstinline{erase(iterator pos);} \lstinline{pop_back();} etc.
  \item
    Stepping through the elements
    \begin{lstlisting}
at(int index) // Random access value.
operator[](int index); // Random access reference
front(); // the first element
back();  // the last element
\end{lstlisting}
  \end{itemize}
\item
  Vector elements can be processed just as they can in an array.
\end{itemize}
\end{frame}

\begin{frame}{Example Operations on a \texttt{\bfseries vector}}

  \centerline{%
  \begin{tikzpicture}[
    node of array/.append style={minimum height=7mm},
    link/.style = {
      -stealth,
      shorten >=1pt,
      shorten <=1pt,
      ultra thick
    },
    array index/.style = {
      font=\small\ttfamily,
    }]
    
    \ArrayLabel{}
    \ArrayList{12, 7, 9, 21, 13}

    \node[block] (declaration) [fit=(el12) (el13)]{};
    \node[anchor=south] at ([yshift=1.1em] declaration.north) {\lstinline|std::vector<int> vec\{12, 7, 9, 21, 13\};|};


    \onslide<2->{
    \begin{scope}[xshift=-3.5cm]
      \ArrayLabel[-2.5]{}
      \ArrayList{12, 7, 9, 21}


      \node[node of array, anchor=west, name=val, xshift=0.75cm, yshift=0.25cm, rotate=35] at ([xshift=-2pt] prev.east) {13};
      \node[block] ({el13}) [fit=(val)]{};

      \draw[link] (el21) to[bend right=10] (el13);
      

      \node[block] (pop) [fit=(el12) (el13)]{};

      \node[anchor=south] at ([yshift=1pt] pop.north) {\lstinline|vec.pop_back();|};
    \end{scope}}

    \onslide<3->{
    \begin{scope}[xshift=3.5cm]
      \ArrayLabel[-2.5]{}
      \ArrayList{12, 7, 9, 21}


      \node[node of array, anchor=west, name=val, xshift=0.75cm, yshift=0.25cm, rotate=35] at ([xshift=-2pt] prev.east) {15};
      \node[block] ({el15}) [fit=(val)]{};

      \draw[link] (el15) to[bend left=10] (el21);
      

      \node[block] (push) [fit=(el12) (el15)]{};

      \node[anchor=south] at ([yshift=1pt] push.north) {\lstinline|vec.push_back(15);|};
    \end{scope}}


    \onslide<4->{
    \ArrayLabel[-4.5]{}
    \ArrayNode{12}
    \ArrayElLabel[val]{vec.begin()}
    \ArrayList{7, 9}
    \ArrayNode{21}
    \ArrayElLabel[val]{vec[3]}
    \ArrayNode{15}}

  \end{tikzpicture}%
  }

\end{frame}

\begin{frame}{Capacity/Size of a \texttt{\bfseries vector}}
\renewcommand{\arraystretch}{1.2}
\begin{tabularx}{1\textwidth}{l X}
\lstinline{empty()} &
  returns \lstinline{true} if the container is empty and \lstinline{false} otherwise \\
\lstinline{size()} &
  returns the number of elements\\
\lstinline{max_size()} &
  returns the maximum possible number of elements\\
\lstinline{reserve()} &
  reserves storage\\
\lstinline{capacity()} &
  returns the number of elements that can be held in currently allocated storage\\
\lstinline{shrink_to_fit()} &
  reduces memory usage by freeing unused memory\\
\end{tabularx}
\end{frame}

\begin{frame}[fragile]{Example use of the \texttt{\bfseries vector} Container}
  \small
\begin{lstlisting}
#include <vector>
#include <iostream>

int main() {
  std::vector<int> v(3);    // a vector of size 3
  v[0] = 23;
  v[1] = 12;
  v[2] = 9;                 // vector full
  v.push_back(17);          // append a new value
  for (auto el : v)         // range-based for loop.
    std::cout << el << ' '; // access to i-th el
  std::cout << std::endl;
  v[2] = v[3];              // random access
}
\end{lstlisting}
\end{frame}

\begin{frame}[fragile]{Example use of the \texttt{\bfseries vector} Container}
  \small
\begin{lstlisting}
#include <vector>
#include <iostream>

int main() {
  std::vector<int> v{ 12, 3, 17, 8 }; // initialize
  while (!v.empty())                  // until empty
  {
      std::cout << v.back() << ' '; // output last el
      v.pop_back();                 // delete last el
  }
  std::cout << std::endl;
}
\end{lstlisting}
\end{frame}

\section{Other Sequence Containers}

\begin{frame}[fragile]{The \texttt{\bfseries deque} (Double-Ended Queue) Container}
\href{https://en.cppreference.com/w/cpp/container/deque}{\lstinline{std::deque}} is another indexed sequence container, which is \\accessed using
\lstinline{#include <deque>}

\begin{itemize}
\item
  Allows fast insertion and deletion at both end ends. Slower insertion in the middle.
\item
  Insertion and deletion at either end of a deque never invalidates pointers or references to the rest of the elements.
\item
  The elements of a deque are not stored contiguously, but random access is $O(1)$.
\end{itemize}
\end{frame}

\begin{frame}[fragile]{The \texttt{\bfseries list}  Container}
\href{https://en.cppreference.com/w/cpp/container/list}{\lstinline{std::list}} is another indexed sequence container, which is accessed using\\
\lstinline{#include <list>}

\begin{itemize}
\item
  Allows constant-time insertion and deletion from anywhere in the container, which never invalidates pointers or references to the rest of the elements.
\item
  Fast random access is not supported.
\item
  Usually implemented as a doubly-linked list.
\item
  Less space efficient.
\end{itemize}
\end{frame}

\section{Iterators}

\begin{frame}{Iterators for Container}
\begin{itemize}
\item
  Iterators are pointer-like entities that are used to access individual elements in a container.
\item
  Often they are used to move sequentially from element to element, a process called \textit{iterating} through a container.
\end{itemize}
\end{frame}


\begin{frame}{Iterators to a \texttt{\bfseries vector} Container}
  \small
  \vspace{1em}
  \begin{tikzpicture}

    % Dynamic Array
    \ArrayLabel{}
    \ArrayNodeVertical{17}
    \ArrayNodeVertical{5}
    \ArrayListVertical{23, 12}
    \ArrayNodeVerticalSpan{3}
    \ArrayNodeVertical{?}
    \ArrayGroupVerticleRight{space}{prev}{Unused}

    % Vector Class Members
    \begin{scope}[node of array/.append style={minimum width=14mm}]
      \ArrayLabel[-2.75][-2.75]{}
      
      \ArrayNodeStack[CSUcyan][capacity]{8}
      \node[block] (longestTitle) [fit=(index)]{};

      \ArrayNodeStack[CSUcyan][size]{4}
      
      \ArrayNodeStack[CSUgold][pArray]{}
      \fill (prev) circle(3pt);
      \coordinate (ptrCircle) at (prev);
      \draw[link] (ptrCircle) to[bend left] (el17);

    \end{scope}
    \begin{scope}[on background layer]
      \node[draw,fill=CSUslate, very thick,fit=(el8) (el4) (el) (longestTitle)] (stackMembers) {};
    \end{scope}
    \node[anchor=south, font=\ttfamily] at ([yshift=0pt] stackMembers.north) {\lstinline|std::vector<int> vec;|};

    % Iterator Members
    \onslide<2->{
      \node[pointer, name=ptr2] at ([yshift=-8.5mm] 2,-1) {};%
      \node[outer block] (last) [fit=(ptr2)]{};
      \node[label of pointer, above = of last.north west, name=head, anchor=south west, align=left, fill=codeBackground] (head) {\lstinline|std::vector<int>::iterator|\\
      \texttt{~~~~}\lstinline|iter = vec.begin() + 1;|};%
      \coordinate (ptr) at (ptr2);%
      \coordinate (headPtr) at (ptr2);%
      \fill (ptr) circle(3pt);%
      \draw[link] (headPtr) to[bend right=8] (el5);
    }

    \onslide<3->{
      \node[below =3pt of last.south west, anchor=north west, align=left, xshift=-3pt, execute at begin node=\setlength{\baselineskip}{1.4em}] (description) {The iterator corresponds to the \\class \lstinline|std::vector<int>| is of the \\type \lstinline|std::vector<int>::iterator|.};
    }
  \end{tikzpicture}
\end{frame}

\begin{frame}[fragile]{Iterators to a \texttt{\bfseries vector} Container}
  \small
\begin{lstlisting}
std::vector<int> vec { 17, 5, 23, 12, 9, 11, 13, 9 };

// Declares an iterator into a int vector container
std::vector<int>::iterator intVecIter = vec.begin();

// Advance the iterator to the second element.
++intVecIter;

// Access the second element and change its value.
*intVecIter = -3;

// Second iterator positioned on the last element.
std::vector<int>::iterator intVecIter2 = vec.end();
--intVecIter2;
\end{lstlisting}
\end{frame}

\begin{frame}[fragile]{\texttt{\bfseries begin} and \texttt{\bfseries end} Iterator Functions}
  \small
\begin{itemize}
\item
  \lstinline{begin()}: returns an iterator to the beginning of the container.
\item
  \lstinline{end()}: returns an iterator to just past the last value.
\end{itemize}
\centerline{%
  \begin{tikzpicture}[node of array/.append style = {minimum height=7mm}]
    \ArrayLabel{\texttt{vec}}
    \ArrayNode[]{17}
    \ArrayElLabel{\lstinline|begin()|}
    \ArrayList{4, 23, 12, 9, 7, 11, 13}
    %\ArrayNodeSpan{5}
    %\ArrayList{,,,,,,}
    \begin{scope}[node of array/.append style={densely dotted}]
      \ArrayNode[none]{}
      \node[anchor=east, inner sep=1pt] (txtLabel) at ([shift={(-0.25,1)}]val.north) {Past the last element};
      \draw[->, very thick] (txtLabel.east) to[bend left=50] (val);
    \end{scope}
    \ArrayElLabel{\lstinline|end()|}
    %\node Past-the-last element
  \end{tikzpicture}
}

\begin{lstlisting}
std::vector<int> vec { 17, 4, 23, 12, 9, 7, 11, 13 };

std::vector<int>::iterator intVecIter1 = vec.begin();
std::vector<int>::iterator intVecIter2 = vec.end();
--intVecIter2; // move to last element
\end{lstlisting}
\end{frame}


\begin{frame}[fragile]{Iterator Example}
  \small
\begin{lstlisting}
template <typename Type>
Type maxValue(const vector<Type> &vec) {
  vector<int>::const_iterator iter = vec.cbegin();
  Type max = *iter;
  for (++iter; iter != vec.cend(); ++iter) {
    if (*iter > max)
      max = *iter;
  }
  return max;
}
\end{lstlisting}

\begin{lstlisting}
std::vector<std::string> words { "The", "fun", "has",
  "just", "begun" };
cout << "max of words = " << maxValue(words);
\end{lstlisting}
\end{frame}

\begin{frame}[fragile]{Types of Iterators}
Every algorithm requires an iterator with a certain capability.
\begin{enumerate}[<+->]
\item
  \emph{Input iterators} have read access; step forward element-by-element.
\item
  \emph{Output iterators} have write access; step forward element-by-element.
\item
  \emph{Forward iterators} have the functionality of input and most of the functionality of output iterators.
\item
  \emph{Bidirectional iterators} are forward iterators that can also go backward.
\item
  \emph{Random access iterators} are bidirectional iterators that can randomly process the elements of a container.
\item As of C++20, there are additional iterators based on \href{https://en.cppreference.com/w/cpp/language/constraints}{concepts}, which differ from C++17.
\end{enumerate}

\href{https://en.cppreference.com/w/cpp/iterator}{https://en.cppreference.com/w/cpp/iterator}
\end{frame}

\begin{frame}[fragile]{Iterator Operations}
\begin{center}\large
\begin{tabular} {l c c c c}
 & \multicolumn{4}{c}{\bfseries Operations}\\
\multicolumn{1}{c}{\bfseries Type} & {Access} & {Read} & {Write} & {Iterate}\\
Input         &               & \lstinline$x = *i$ &               &  \lstinline$++$\\
Output        &               & \lstinline$x = *i$ &  \lstinline$*i = x$ &  \lstinline$++$\\
Forward       &               & \lstinline$x = *i$ &  \lstinline$*i = x$ &  \lstinline$++$\\
Bidirectional &               & \lstinline$x = *i$ &  \lstinline$*i = x$ &  \lstinline$++$, \lstinline$--$\\
Random Access &\lstinline$[ ]$ & \lstinline$x = *i$ &  \lstinline$*i = x$ &  \lstinline$++$, \lstinline$--$\\
\end{tabular}
\end{center}
\end{frame}

\begin{frame}[fragile]{Helper Iterator Operators}
  The \lstinline|<iterator>| library includes many overloaded operators and functions to use with iterators. Some include:
  \begin{itemize}
    \item \href{https://en.cppreference.com/w/cpp/iterator/advance}{advance} -- advances an iterator by a given distance
    \item \href{https://en.cppreference.com/w/cpp/iterator/next}{next} -- increment an iterator
    \item \href{https://en.cppreference.com/w/cpp/iterator/prev}{prev} -- decrement an iterator 
  \end{itemize} 

\end{frame}

\begin{frame}{Benefits of Iterators}
\begin{itemize}
\item<2->
  \emph{Convenience in Programming}:
  \begin{itemize}
  \item
    Provide common usage for any container (the internal structure of a container does not matter).
  \item
    Provide read and/or write access to each element's values.
  \end{itemize}
\item<3->
  \emph{Reusable Code}:\\
  Iterator algorithms are not dependent on the container type.
\item<4->
  \emph{Dynamic Processing}:\\
  Easily dynamically allocate as well as delete the memory.
\end{itemize}
\end{frame}

\section{Associative Containers}

\begin{frame}{Ordered Associative Containers}

Store elements in sorted order according to some ordering criteria\\(often using a type of binary tree).\\[1em]

\begin{tabular}{l l l}
  \textbf{\large Associative Containers} & \textbf{Header File} & \textbf{Iterator Support}\\
  \href{https://en.cppreference.com/w/cpp/container/map}{map}
    & \lstinline{<map>} & Bidirectional\\
  \href{https://en.cppreference.com/w/cpp/container/multima}{multimap}
    & \lstinline{<map>} & Bidirectional\\
  \href{https://en.cppreference.com/w/cpp/container/set}{set}
    & \lstinline{<set>} & Bidirectional\\
  \href{https://en.cppreference.com/w/cpp/container/multiset}{multiset}
    & \lstinline{<set>} & Bidirectional\\
\end{tabular}

\begin{itemize}
\item
  \lstinline{set} stores a set of \emph{unique} elements, while \lstinline{map} stores key-value pairs with \emph{unique} keys.
\item
  \lstinline{multiset} and \lstinline{multimap} allow duplicates.
\end{itemize}

\end{frame}

\begin{frame}[fragile]{Sets Example}
\small
\begin{lstlisting}[basicstyle=\footnotesize\ttfamily]
#include <set>
int main() {
  std::set<std::string, std::less<std::string> > nameSet{ "Ole",
       "Hedvig", "Juan", "Lars", "Guido", "Ann"};

  nameSet.insert("Patric"); // insert another name
  nameSet.erase("Juan"); // removes an element

  // Find the matching name in set.
  std::string searchName;
  std::cin >> searchName;
  std::set<std::string, std::less<string> >::iterator iter =
    nameSet.find(searchName);

  cout << searchName << " is";
  if (iter == nameSet.end()) // check if found
     std::cout << " NOT";
  std::cout << " in the set!\n";
}
\end{lstlisting}
\end{frame}


\begin{frame}[fragile]{Sets Example}
\small
\begin{lstlisting}[basicstyle=\footnotesize\ttfamily]
#include <set>

int main() {
  std::set<std::string, std::less<std::string> > nameSet {
    "Ole", "Hedvig", "Juan", "Lars", "Guido", "Patric",
    "Maria", "Ann" };

  // Set iterator to lower start value "K"
  std::set<std::string, std::less<std::string> >::iterator
    iter = nameSet.lower_bound("K");

  // Display: Lars, Maria, Ole, Patric
  while (iter != nameSet.upper_bound("Q"))
    cout << *iter++ << ", ";
  std::cout << std::endl;
}
\end{lstlisting}
\end{frame}

\begin{frame}[fragile]{Maps Example}
\begin{lstlisting}[basicstyle=\footnotesize\ttfamily]
#include <map>
#include <algorithm>
int main() {
  std::map<std::string, int> contacts {
    {"Ole", 75643}, {"Hedvig", 83268}, {"Juan", 97353},
    {"Lars", 87353}, {"Guido", 19988}, {"Patric", 76455},
    {"Maria", 77443}, {"Ann", 12221}};

  std::cout << "Lars phone number is " << contacts["Lars"]
    << '\n';

  for(auto iter = contacts.begin(); iter != contacts.end(); iter++)
    std::cout << iter->first << " : " << iter->second << '\n';
}
\end{lstlisting}
\end{frame}


\section{Adapters}

\begin{frame}{Container Adapters}

Containers to accommodate special situations.\\[1em]

\begin{tabular}{l l l}
\toprule
  \textbf{\large Adapters} & \textbf{Header File} & \textbf{Iterator Support}\\
\midrule
  \href{https://en.cppreference.com/w/cpp/container/stack}{stack}
    & \lstinline{<stack>} & No iterator support\\
  \href{https://en.cppreference.com/w/cpp/container/queue}{queue}
    & \lstinline{<queue>} & No iterator support\\
  \href{https://en.cppreference.com/w/cpp/container/priority_queue}{priority queue}
    & \lstinline{<queue>} & No iterator support\\
\bottomrule
\end{tabular}

\begin{itemize}
\item
  Do not support any type of iterator.
\end{itemize}

\end{frame}

\begin{frame}{The \texttt{\textbf{std::span}} Container View}
  \begin{tikzpicture}

    % Dynamic Array
    \ArrayLabel{\lstinline|int array[12];|}
    \ArrayNodeVertical{2}
    \ArrayNodeVertical{3}
    \ArrayNodeVertical{5}
    \ArrayNodeVertical{7}
    \ArrayNodeVerticalSpan{7}
    \ArrayNodeVertical{37}
    %\ArrayGroupVerticleRight{space}{prev}{Unused}

    % Span 1 Class Members
    \begin{scope}[node of array/.append style={minimum width=14mm}]
      \ArrayLabel[-2.5][8]{}
      
      %\ArrayNodeStack[CSUcyan][capacity]{8}
      
      \ArrayNodeStack[CSUcyan][size]{12}
      
      \ArrayNodeStack[CSUgold][pointer]{}
      \node[block] (longestTitle) [fit=(index)]{};
      \fill (prev) circle(3pt);
      \coordinate (ptrCircle) at (prev);
      \draw[link] (ptrCircle) to[bend right=8] (el2);

    \end{scope}
    \begin{scope}[on background layer]
      \node[draw,fill=CSUslate, very thick,fit=(el12) (el) (longestTitle)] (stackMembers) {};
    \end{scope}
    \node[anchor=south, font=\ttfamily] at ([yshift=0pt] stackMembers.north) {\lstinline|std::span<int> arrSpan\{array, 12\};|};


    % Span 2 Class Members
    \onslide<2->{
    \begin{scope}[node of array/.append style={minimum width=14mm}]
      \ArrayLabel[-6][7]{}
      
      %\ArrayNodeStack[CSUcyan][capacity]{8}
      
      \ArrayNodeStack[CSUcyan][size]{9}
      
      \ArrayNodeStack[CSUgold][pointer]{}
      \node[block] (longestTitle) [fit=(index)]{};
      \fill (prev) circle(3pt);
      \coordinate (ptrCircle) at (prev);
      \draw[link] (ptrCircle) to[bend right=8] (el7);

      \end{scope}
      \begin{scope}[on background layer]
        \onslide<2->{
          \node[draw,fill=CSUslate, very thick,fit=(el9) (el) (longestTitle)] (span2Members) {};
        }
      \end{scope}
      \node[anchor=south, font=\ttfamily] at ([yshift=0pt] span2Members.north) {\lstinline|std::span<int> part\{array + 3, 9\};|};
    }
    % % Iterator Members
    % \onslide<2->{
    %   \node[pointer, name=ptr2] at ([yshift=-8.5mm] 2,-0.5) {};%
    %   \node[outer block] (last) [fit=(ptr2)]{};
    %   \node[label of pointer, above = of last.north west, name=head, anchor=south west] (head) {\lstinline|std::vector<int>::iterator|};%
    %   \coordinate (ptr) at (ptr2);%
    %   \coordinate (headPtr) at (ptr2);%
    %   \fill (ptr) circle(3pt);%
    %   \draw[link] (headPtr) to[bend right=8] (el5);
    % }

    % \onslide<3->{
    %   \node[below =3pt of last.south west, anchor=north west, align=left, xshift=-3pt, execute at begin node=\setlength{\baselineskip}{1.4em}] (description) {The iterator corresponds to the \\class \lstinline|std::vector<int>| is of the \\type \lstinline|std::vector<int>::iterator|.};
    % }
  \end{tikzpicture}
\end{frame}

\section{Algorithms}

\begin{frame}{Algorithms in the STL}
Many algorithms that work on any STL container are \\
found in the \lstinline{<algorithm>} library.
\begin{itemize}
\item
  Copy
\item
  Partition
\item
  Sort
\item
  Search
\item
  \ldots{}
\end{itemize}

These algorithms are provided generically to support many container types.
\end{frame}


\begin{frame}{STL Algorithm Classifications}

\begin{itemize}
\item
  \emph{Non-modifying algorithms} do not modify the elements of the container
\item
  \emph{Modifying algorithms} modify the elements of the container
\item
  \emph{Numeric algorithms} perform numeric calculations on elements of a container.
\item
  \emph{Heap algorithms} are Heapsort algorithm that sorts contiguous data.
\end{itemize}
\end{frame}

\begin{frame}{Non-Modifying Sequence Algorithms}
\footnotesize\renewcommand{\ULdepth}{3pt}\renewcommand{\arraystretch}{1.2}
\begin{tabularx}{\textwidth}{r X}

\href{https://en.cppreference.com/w/cpp/algorithm/all_any_none_of}{\lstinline{all\_of, any\_of, none\_of}}
 &    checks if a predicate is true for all, any or none of the elements in a range \\
\href{https://en.cppreference.com/w/cpp/algorithm/for_each}{\lstinline{for\_each}}
 &    applies a function to a range of elements \\
\href{https://en.cppreference.com/w/cpp/algorithm/for_each_n}{\lstinline{for\_each\_n}}
 &    applies a function object to the first n elements of a sequence  \\
\href{https://en.cppreference.com/w/cpp/algorithm/count}{\lstinline{count, count\_if}}
 &    returns the number of elements satisfying specific criteria  \\
\href{https://en.cppreference.com/w/cpp/algorithm/mismatch}{\lstinline{mismatch}}
 &    finds the first position where two ranges differ \\
\href{https://en.cppreference.com/w/cpp/algorithm/find}{\lstinline{find, find\_if, find\_if\_not}}
 &    finds the first element satisfying specific criteria  \\
\href{https://en.cppreference.com/w/cpp/algorithm/find_end}{\lstinline{find\_end}}
 &    finds the last sequence of elements in a certain range \\
\href{https://en.cppreference.com/w/cpp/algorithm/find_first_of}{\lstinline{find\_first\_of}}
 &    searches for any one of a set of elements  \\
\href{https://en.cppreference.com/w/cpp/algorithm/adjacent_find}{\lstinline{adjacent\_find}}
 &    finds the first two adjacent items that are equal (or satisfy a given predicate)  \\
\href{https://en.cppreference.com/w/cpp/algorithm/search}{\lstinline{search}}
 &    searches for a range of elements \\
\href{https://en.cppreference.com/w/cpp/algorithm/search_n}{\lstinline{search\_n}}
 &    searches a range for a number of consecutive\\[-0.2em]
 &    copies of an element \\
\end{tabularx}
\end{frame}


\begin{frame}[fragile]{Example: The \texttt{for\_each()} Algorithm}
  \small
\begin{lstlisting}
#include <vector>
#include <algorithm>
#include <iostream>
void show(int n) {
  std::cout << n << ", ";
}

int main() {
  // initialize vector with a list of values
  std::vector<int> v{ 12, 3, 17, 8 };

  // apply function show to each element of vector v
  std::for_each (v.begin(), v.end(), show);
}
\end{lstlisting}
\end{frame}

\begin{frame}[fragile]{Example: The \texttt{find()} Algorithm}
  
\begin{lstlisting}[basicstyle=\small\ttfamily]
#include <vector>
#include <algorithm>
#include <iostream>
int main() {
  std::vector<int> v{ 12, 3, 17, 8, 34, 56, 9 };
  std::vector<int>::iterator iter;
  int key;

  std::cout << "Enter value: ";
  std::cin >> key;
  iter = std::find(v.begin(), v.end(), key);
  std::cout << "Element " << key << " is ";
  if (iter == v.end()) // value was NOT found
     std::cout << "not";
  std::cout << " in vector v.\n";
}
\end{lstlisting}
\end{frame}

\begin{frame}[fragile]{Example: The \texttt{find\_if()} Algorithm}
\begin{lstlisting}[basicstyle=\small\ttfamily]
#include <vector>
#include <algorithm>
#include <iostream>
bool inRange(int n) { return (n > 21) && (n < 36); };

int main() {
  std::vector<int> v{ 12, 3, 17, 8, 34, 56, 9 };
  std::vector<int>::iterator iter;
  iter = std::find_if(v.begin(), v.end(), inRange);

  // finds an element in v for which inRange is true
  if (iter != v.end()) // found the element
    std::cout << "found " << *iter << '\n';
  else
    std::cout << "not found\n";
}
\end{lstlisting}
\end{frame}

\begin{frame}[fragile]{Example: The \texttt{count\_if()} Algorithm}
\begin{lstlisting}[basicstyle=\small\ttfamily]
#include <vector>
#include <algorithm>
#include <iostream>
bool inRange(int n) { return (n > 14) && (n < 36); };

int main() {
  std::vector<int> v{ 12, 3, 17, 8, 34, 56, 9 };
  std::vector<int>::iterator iter;

  // counts element in v for which inRange is true
  const int N = std::count_if(v.begin(),v.end(), inRange);
  std::cout << "found " << N << " values between 14 and 36.\n";
}
\end{lstlisting}
\end{frame}

\begin{frame}{Minimum/Maximum Algorithms}
\small\renewcommand{\ULdepth}{3pt}\renewcommand{\arraystretch}{1.3}
\begin{tabularx}{\textwidth}{r X}
\href{https://en.cppreference.com/w/cpp/algorithm/max}{\lstinline{max}} &
  returns the greater of the given values\\
\href{https://en.cppreference.com/w/cpp/algorithm/max_element}{\lstinline{max_element}} &
  returns the largest element in a range\\
\href{https://en.cppreference.com/w/cpp/algorithm/min}{\lstinline{min}} &
  returns the smaller of the given values\\
\href{https://en.cppreference.com/w/cpp/algorithm/min_element}{\lstinline{min_element}} &
  returns the smallest element in a range\\
\href{https://en.cppreference.com/w/cpp/algorithm/minmax}{\lstinline{minmax}} &
  returns the smaller and larger of two elements\\
\href{https://en.cppreference.com/w/cpp/algorithm/minmax_element}{\lstinline{minmax_element}} &
  returns the smallest and the largest elements in a range\\
\href{https://en.cppreference.com/w/cpp/algorithm/clamp}{\lstinline{clamp}} &
  clamps a value between a pair of boundary values\\
\end{tabularx}
\end{frame}

\begin{frame}{Set Algorithms (on Sorted Ranges)}
\small\renewcommand{\ULdepth}{3pt}\renewcommand{\arraystretch}{1.3}
\begin{tabularx}{\textwidth}{r X}
\href{https://en.cppreference.com/w/cpp/algorithm/includes}{\lstinline{includes}} &
  returns true if one sequence is a subsequence of another\\
\href{https://en.cppreference.com/w/cpp/algorithm/set_difference}{\lstinline{set_difference}} &
  computes the difference between two sets\\
\href{https://en.cppreference.com/w/cpp/algorithm/set_intersection}{\lstinline{set_intersection}} &
  computes the intersection of two sets\\
\href{https://en.cppreference.com/w/cpp/algorithm/set_symmetric_difference}{\lstinline{set_symmetric_difference}} &
  computes the symmetric difference between two sets\\
\href{https://en.cppreference.com/w/cpp/algorithm/set_union}{\lstinline{set_union}} &
  computes the union of two sets\\
\end{tabularx}
\end{frame}

\begin{frame}{Modifying Sequence Algorithms}
\footnotesize\renewcommand{\ULdepth}{3pt}\renewcommand{\arraystretch}{1.3}
\begin{tabularx}{\textwidth}{r X}
\href{https://en.cppreference.com/w/cpp/algorithm/copy}{\lstinline{copy, copy_if}} &
  copies a range of elements to a new location\\
\href{https://en.cppreference.com/w/cpp/algorithm/copy_n}{\lstinline{copy_n}} &
  copies a number of elements to a new location\\
\href{https://en.cppreference.com/w/cpp/algorithm/copy_backward}{\lstinline{copy_backward}} &
  copies a range of elements in backward order\\
\href{https://en.cppreference.com/w/cpp/algorithm/move}{\lstinline{move}} &
  moves a range of elements to a new location\\
\href{https://en.cppreference.com/w/cpp/algorithm/move_backward}{\lstinline{move_backward}} &
  moves a range of elements to a new location in backward order\\
\href{https://en.cppreference.com/w/cpp/algorithm/fill}{\lstinline{fill}} &
  copy-assigns the given value to every element in a range\\
\href{https://en.cppreference.com/w/cpp/algorithm/fill_n}{\lstinline{fill_n}} &
  copy-assigns the given value to N elements in a range\\
\href{https://en.cppreference.com/w/cpp/algorithm/transform}{\lstinline{transform}} &
  applies a function to a range of elements, storing results in a destination range\\
\href{https://en.cppreference.com/w/cpp/algorithm/generate}{\lstinline{generate}} &
  assigns the results of successive function calls to every element in a range\\
\href{https://en.cppreference.com/w/cpp/algorithm/generate_n}{\lstinline{generate_n}} &
  assigns the results of successive function calls to N elements in a range\\
\href{https://en.cppreference.com/w/cpp/algorithm/remove}{\lstinline{remove, remove_if}} &
  removes elements satisfying specific criteria\\

\end{tabularx}
\end{frame}

\begin{frame}[fragile]{Example: The \texttt{copy()} Algorithm}

The copy algorithm overwrites the existing contents when using regular iterators.

\begin{lstlisting}[basicstyle=\small\ttfamily]
#include <list>
#include <algorithm>

int main() {
  std::list<int> odds{ 1, 3, 5, 7, 9 };
  std::list<int> result{ 2, 4, 6, 8, 10 };

  // copy contents of odds to result overwriting the elements in result
  std::copy(odds.begin(), odds.end(), result.begin());
  // result = { 1, 3, 5, 7, 9 }
}
\end{lstlisting}
\end{frame}

\begin{frame}[fragile]{Example: The \texttt{copy()} Algorithm}
\small
With insert operators, you can modify the behavior of the copy algorithm.\\[0.5em]

\begin{tabular}{l l}
\href{https://en.cppreference.com/w/cpp/iterator/back_inserter}{\lstinline{back_inserter}} & inserts new elements at the end\\
\href{https://en.cppreference.com/w/cpp/iterator/front_inserter}{\lstinline{front_inserter}} & inserts new elements at the beginning\\
\href{https://en.cppreference.com/w/cpp/iterator/inserter}{\lstinline{inserter}} & inserts new elements at a specified location\\
\end{tabular}\\[1em]

\begin{lstlisting}[basicstyle=\small\ttfamily]
std::list<int> l1{ 1, 3, 5, 7, 9 };
std::list<int> l2{ 2, 4, 6, 8, 10 };
\end{lstlisting}

\begin{lstlisting}[basicstyle=\footnotesize\ttfamily]
std::copy(l1.begin(), l1.end(), std::back_inserter(l2));
// l2 = { 2, 4, 6, 8, 10, 1, 3, 5, 7, 9  }
\end{lstlisting}

\begin{lstlisting}[basicstyle=\footnotesize\ttfamily]
std::copy(l1.begin(), l1.end(), std::front_inserter(l2));
// l2 = { 9, 7, 5, 3, 1, 2, 4, 6, 8, 10 }
\end{lstlisting}

\begin{lstlisting}[basicstyle=\footnotesize\ttfamily]
std::copy(l1.begin(), l1.end(), std::inserter(l2, l2.begin()));
// l2 = { 1, 3, 5, 7, 9, 2, 4, 6, 8, 10 }
\end{lstlisting}
\end{frame}

\begin{frame}[fragile]{Example: \texttt{copy()} with Output Iterator}

\begin{lstlisting}[basicstyle=\small\ttfamily]
#include <iostream>
#include <iterator>
#include <vector>
#include <algorithm>

int main() {
  std::vector<int> vector {10, 20, 30, 40, 50, 60 70};
  std::ostream_iterator<int> outIter{std::cout, ", "};

  std::copy(vector.begin(), vector.end(), outIter);
}
\end{lstlisting}
\end{frame}

\begin{frame}{Modifying Sequence Algorithms}
\footnotesize\renewcommand{\ULdepth}{3pt}\renewcommand{\arraystretch}{1.3}
\begin{tabularx}{\textwidth}{r X}
\href{https://en.cppreference.com/w/cpp/algorithm/remove_copy}{\lstinline{remove_copy, remove_copy_if}} &
  copies a range of elements omitting those that satisfy specific criteria\\
\href{https://en.cppreference.com/w/cpp/algorithm/replace}{\lstinline{replace, replace_if}} &
  replaces all values satisfying specific criteria with another value\\
\href{https://en.cppreference.com/w/cpp/algorithm/replace_copy}{\lstinline{replace_copy, replace_copy_if}} &
  copies a range, replacing elements satisfying specific criteria with another value\\
\href{https://en.cppreference.com/w/cpp/algorithm/swap}{\lstinline{swap}} &
  swaps the values of two objects\\
\href{https://en.cppreference.com/w/cpp/algorithm/swap_ranges}{\lstinline{swap_ranges}} &
  swaps two ranges of elements\\
\href{https://en.cppreference.com/w/cpp/algorithm/iter_swap}{\lstinline{iter_swap}} &
  swaps the elements pointed to by two iterators\\
\href{https://en.cppreference.com/w/cpp/algorithm/reverse}{\lstinline{reverse}} &
  reverses the order of elements in a range\\
\href{https://en.cppreference.com/w/cpp/algorithm/reverse_copy}{\lstinline{reverse_copy}} &
  creates a copy of a range that is reversed\\
\href{https://en.cppreference.com/w/cpp/algorithm/rotate}{\lstinline{rotate}} &
  rotates the order of elements in a range\\
\href{https://en.cppreference.com/w/cpp/algorithm/rotate_copy}{\lstinline{rotate_copy}} &
  copies and rotate a range of elements\\
\end{tabularx}
\end{frame}

\begin{frame}[fragile]{Example: The \texttt{sort()} and \texttt{merge()} Algorithms}

\begin{lstlisting}[basicstyle=\small\ttfamily]
#include <list>
#include <algorithm>
int main() {
  std::list<int>  l1 { 6, 4, 9, 1, 7 }
  std::list<int>  l2 { 4, 2, 1, 3, 8 };
  l1.sort(); // l1 = {1, 4, 6, 7, 9}
  l2.sort(); // l2 = {1, 2, 3, 4, 8 }
  l1.merge(l2);  // merges l2 into l1
  // l1 = { 1, 1, 2, 3, 4, 4, 6, 7, 8, 9},  l2 = {}
}
\end{lstlisting}
\end{frame}

\begin{frame}{Numeric Algorithms in \href{https://en.cppreference.com/w/cpp/header/numeric}{\texttt{<numeric>}}}
\footnotesize\renewcommand{\ULdepth}{3pt}\renewcommand{\arraystretch}{1.3}
\begin{tabularx}{\textwidth}{r X}
\href{https://en.cppreference.com/w/cpp/algorithm/iota}{\lstinline{iota}} &
  fills a range with successive increments of the starting value\\
\href{https://en.cppreference.com/w/cpp/algorithm/accumulate}{\lstinline{accumulate}} &
  sums up a range of elements\\
\href{https://en.cppreference.com/w/cpp/algorithm/inner_product}{\lstinline{inner_product}} &
  computes the inner product of two ranges of elements\\
\href{https://en.cppreference.com/w/cpp/algorithm/adjacent_difference}{\lstinline{adjacent_difference}} &
  computes the differences between adjacent elements in a range\\
\href{https://en.cppreference.com/w/cpp/algorithm/partial_sum}{\lstinline{partial_sum}} &
  computes the partial sum of a range of elements\\
\href{https://en.cppreference.com/w/cpp/algorithm/reduce}{\lstinline{reduce}} &
  similar to std::accumulate, except out of order\\
\href{https://en.cppreference.com/w/cpp/algorithm/exclusive_scan}{\lstinline{exclusive_scan}} &
  excludes the ith input element from the ith sum\\
\href{https://en.cppreference.com/w/cpp/algorithm/inclusive_scan}{\lstinline{inclusive_scan}} &
  includes the ith input element in the ith sum\\
\href{https://en.cppreference.com/w/cpp/algorithm/transform_reduce}{\lstinline{transform_reduce}} &
  applies an invocable, then reduces out of order\\
\href{https://en.cppreference.com/w/cpp/algorithm/transform_exclusive_scan}{\lstinline{transform_exclusive_scan}} &
  applies an invocable, then calculates exclusive scan\\
\href{https://en.cppreference.com/w/cpp/algorithm/transform_inclusive_scan}{\lstinline{transform_inclusive_scan}} &
  applies an invocable, then calculates inclusive scan\\
\end{tabularx}
\end{frame}

\begin{frame}{Heap Algorithms}

Built-in heapsort.\\
(You may not implement the Heap in your lab using these functions,\\ but you may use them for testing.)

\footnotesize\renewcommand{\ULdepth}{3pt}\renewcommand{\arraystretch}{1.3}
\begin{tabularx}{\textwidth}{r X}
\href{https://en.cppreference.com/w/cpp/algorithm/is_heap}{\lstinline{is_heap}} &
  checks if the given range is a max heap\\
\href{https://en.cppreference.com/w/cpp/algorithm/is_heap_until}{\lstinline{is_heap_until}} &
  finds the largest subrange that is a max heap\\
\href{https://en.cppreference.com/w/cpp/algorithm/make_heap}{\lstinline{make_heap}} &
  creates a max heap out of a range of elements\\
\href{https://en.cppreference.com/w/cpp/algorithm/push_heap}{\lstinline{push_heap}} &
  adds an element to a max heap\\
\href{https://en.cppreference.com/w/cpp/algorithm/pop_heap}{\lstinline{pop_heap}} &
  removes the largest element from a max heap\\
\href{https://en.cppreference.com/w/cpp/algorithm/sort_heap}{\lstinline{sort_heap}} &
  turns a max heap into a range of elements sorted in ascending order\\
\end{tabularx}
\end{frame}

\section{Function Objects}
\begin{frame}[fragile]{Function Objects}
\small
\begin{itemize}
\item
  Some algorithms like sort, merge, and accumulate can take a function object as an argument.
\item
  A function object is an instance of a class that overloads \lstinline{operator()}.
\item
  It is also possible to use user-written functions in place of predefined function objects.
\end{itemize}
\small
\begin{lstlisting}[basicstyle=\small\ttfamily]
#include <list>
#include <functional>
int main() {
  std::list<int> lst { 6, 4, 9, 1, 7 };
  lst.sort(std::greater<int>{}); // uses function object
  // for sorting in reverse order lst = { 9, 7, 6, 4, 1 }
}
\end{lstlisting}
In C++, Function Objects are often called \emph{functors}, but that term has alternate means.
\end{frame}


\begin{frame}[fragile]{Function Objects}
\small
The accumulate algorithm accumulates values over a range (e.g., the sum).

\begin{lstlisting}[basicstyle=\small\ttfamily]
#include <list>
#include <functional>
#include <numeric>
int main() {
  std::list<int> lst { 6, 4, 9, 1, 7 };
  int sum, fac;
  sum = std::accumulate(lst.begin(), lst.end(), 0,
    std::plus<int>{});
  sum = std::accumulate(lst.begin(), lst.end(), 0);
    // plus by default

  fac = std::accumulate(lst.begin(), lst.end(), 1,
    std::times<int>{});
}
\end{lstlisting}
\end{frame}

\begin{frame}[fragile]{User Defined Function Objects}

\begin{lstlisting}[basicstyle=\footnotesize\ttfamily]
template <typename Type>
struct SquaredSum  // user-defined function object
{
    Type operator()(Type n1, Type n2) { return n1 + n2 * n2; }
};
\end{lstlisting}
\begin{lstlisting}[basicstyle=\footnotesize\ttfamily]
#include <complex>
#include <numeric>
#include <vector>
\end{lstlisting}
\begin{lstlisting}[basicstyle=\footnotesize\ttfamily]
std::vector<std::complex<int>> vc {{2, 3}, {1, 5}, {-2, 4}};

// computes the sum of squares of a vector of complex numbers
std::complex sum_vc = std::accumulate(vc.begin(), vc.end(),
          std::complex(0, 0), SquaredSum<std::complex<int>>{});
\end{lstlisting}
\end{frame}


\begin{frame}[fragile]{User Defined Function Objects}
Unlike functions, function objects permit custom parameterization and statefullness.

\begin{lstlisting}[basicstyle=\footnotesize\ttfamily]
std::vector<double> vals1 {2, 3, 1, 5, -2, 4};
std::vector<double> vals2 {9, 8, -6, -10, -1, 8, 6};

// Square each value in both vectors and calculate the total.
CalculateAverageOfPowers squared{2}; // Function Object
squared = std::for_each(vals1.begin(), vals1.end(), squared);
squared = std::for_each(vals2.begin(), vals2.end(), squared);

std::cout << ' ' << squared.getAverage() << '\n';
\end{lstlisting}
\end{frame}

\begin{frame}[fragile]{User Defined Function Objects}
\begin{lstlisting}[basicstyle=\footnotesize\ttfamily]
class CalculateAverageOfPowers
{
public:
  CalculateAverageOfPowers(double power) 
    : total{0}, count{0}, exponent{power} {}

  void operator() (double x) {
    total += pow(x, exponent);
    ++count;
  }

  auto getAverage() const { return total / count; }

private:
  double total;
  int count;
  double exponent;
};
\end{lstlisting}
\end{frame}

\section*{Lab 17}
\begin{frame}{Lab 17: The STL and Sorting}
\centering\Large
  Let's take a look at Lab 17.
\end{frame}

\end{document}
